% Generated from locale/tr/content/set-theory/spine/foundation.tex; do not hand-edit.


\olfileid[tr]{sth}{spine}{foundation}

\olsection{Temellendirme}

Yalnızca \emph{neredeyse} tamamladık---henüz \emph{tam} bitirmedik---çünkü
$\ZFminus$ içinde \emph{her} kümenin bir $V_\alpha$'da bulunduğunu, yani
her kümenin bir aşamada oluşturulduğunu güvence altına alan hiçbir şey
yoktur.

Hiyerarşinin dışında kümeler bulunup bulunmadığını matematiksel açıdan
\emph{önemsemememizin} oldukça açık bir anlamı vardır. (Orada herhangi bir
küme varsa onu görmezden gelebiliriz.) Fakat küme \emph{kavramımızı}, her
kümenin bir aşamada oluşturulduğu düşüncesiyle gerekçelendirdik
(\olref[z][story]{sec} içindeki \stageshier{} ilkesine bakın). Bu yüzden
hiyerarşinin dışında kalan kümeler bulunması olasılığını dışlamak isteriz.
Buna göre, her kümenin hiyerarşinin bir yerinde ortaya çıkmasını güvence
altına alan yeni bir aksiyom eklemeliyiz.

$V_\alpha$'lar aşamalarımız olduğuna göre, şu ifadeyi doğrudan aksiyom
olarak eklemeyi düşünebiliriz:

\begin{defish}
\emph{Düzenlilik.} $\forall A \exists \alpha\, A \subseteq V_\alpha$
\end{defish}

Bu bütünüyle makul bir yaklaşım olurdu. Ancak sonraki kesimde açıklanacak
nedenlerle bunun yerine alternatif bir aksiyom kabul edeceğiz:

\begin{axiom}[Temellendirme]
$(\forall A \neq \emptyset)(\exists B \in A)A \cap B = \emptyset$.
\end{axiom}

Biraz çabayla Temellendirme'nin ($\ZFminus$ içinde) Düzenlilik'i
gerektirdiğini gösterebiliriz:
\begin{defn}
Her $A$ kümesi için:
\begin{align*}
	\text{cl}_0(A) &= A,\\
	\text{cl}_{n+1}(A) &= \bigcup \text{cl}_n(A),\\
	\text{trcl}(A) &= \bigcup_{n < \omega} \text{cl}_{n}(A).
\end{align*}
$\text{trcl}(A)$'ya $A$'nın \emph{geçişli kapanışı} deriz.
\end{defn}
\noindent
``Geçişli kapanış'' adı yerindedir:
\begin{prop}\ollabel{subsetoftrcl}
$A \subseteq \trcl{A}$ ve $\trcl{A}$ geçişli bir kümedir.
\end{prop}

\begin{proof}
Açıktır ki $A=\text{cl}_0(A)\subseteq\text{trcl}(A)$ olur. Ayrıca
$x\in b\in\trcl{A}$ ise bir $n$ için $b\in\text{cl}_n(A)$ ve dolayısıyla
$x\in\text{cl}_{n+1}(A)\subseteq\text{trcl}(A)$ olur.
\end{proof}

\begin{lem}\ollabel{lem:TransitiveWellFounded}
$A$ geçişli bir kümeyse $A\subseteq V_\alpha$ olacak biçimde bir $\alpha$
vardır.
\end{lem}

\begin{proof}
\olref[ordinals][opps]{defsupstrict} içindeki ``$\supstrict(X)$'' tanımını
hatırlayarak iki küme tanımlayın:
\begin{align*}
	D  &= \Setabs{x \in A}{\forall \delta\ x \nsubseteq V_\delta}\\
	\alpha &= \supstrict\Setabs{\delta}{(\exists x \in A)
	(x \subseteq V_\delta \land (\forall \gamma \in \delta)x \nsubseteq V_\gamma)}.
\end{align*}
$D=\emptyset$ olduğunu varsayın. $x\in A$ ise o zaman
$x\subseteq V_\delta$ olacak biçimde bir $\delta$ vardır ve sıra sayılarının
iyi sıralanması gereği
$(\forall\gamma\in\delta)x\nsubseteq V_\gamma$ olan en küçük böyle
$\delta$ seçilebilir. Dolayısıyla $\delta\in\alpha$ ve
\olref[spine][Valphabasic]{Valphabasicprops} gereği $x\in V_\alpha$ olur.
Böylece gerektiği gibi $A\subseteq V_\alpha$ elde edilir.

O hâlde $D=\emptyset$ olduğunu göstermek yeterlidir. Olmayana ergiyle aksi
durumu varsayın. Temellendirme gereği
$B\in D\subseteq A$ ve $D\cap B=\emptyset$ olan bir $B$ vardır.
$x\in B$ ise $A$ geçişli olduğundan $x\in A$ olur; ayrıca $x\notin D$
olduğundan $x\subseteq V_\delta$ olacak biçimde bir $\delta$ vardır. Şimdi
\[
	\beta = \supstrict\Setabs{\delta}{(\exists x \in B)(x \subseteq V_\delta \land (\forall \gamma < \delta)x \nsubseteq V_\gamma)}
\]
olsun. Daha önceki gibi $B\subseteq V_\beta$ elde edilir; bu, $B\in D$
iddiasıyla çelişir.
\end{proof}

\begin{thm}\ollabel{zfentailsregularity}
Düzenlilik geçerlidir.
\end{thm}

\begin{proof}
$A$ sabit olsun. \olref{subsetoftrcl} gereği
$A\subseteq\trcl{A}$ ve $\trcl{A}$ geçişlidir. Dolayısıyla
\olref{lem:TransitiveWellFounded} gereği
$A\subseteq\trcl{A}\subseteq V_\alpha$ olacak biçimde bir $\alpha$ vardır.
\end{proof}

Bu sonuçlar $\ZFminus$'ın
$\emph{Temellendirme}\Rightarrow\emph{Düzenlilik}$ koşullusunu
ispatladığını gösterir. \olref[spine][rank]{zfminusregularityfoundation}
içinde $\ZFminus$'ın
$\emph{Düzenlilik}\Rightarrow\emph{Temellendirme}$ koşullusunu
ispatladığını göstereceğiz. Buna göre Temellendirme ile Düzenlilik
$\ZFminus$ altında \emph{denktir}. Fakat bu, $\ZFminus$ verildiğinde
Temellendirme'yi Düzenlilik'e denk olduğunu belirterek
gerekçelendirebileceğimiz anlamına gelir. Düzenlilik'i ise
\stageshier{} temelinde doğrudan gerekçelendirebiliriz.

